\documentclass[12pt]{article}
\usepackage[a4paper,margin=1in]{geometry}
\usepackage[utf8]{inputenc}
\usepackage[T1]{fontenc}
\usepackage{amsmath,amssymb}
\usepackage{textcomp}
\usepackage{parskip}
\begin{document}
\begin{center}
{\Large \textbf{2012 STEP II}}
\end{center}

\subsection*{Page 1}
\begin{flushleft}
\phi\par
\bigskip\textbf{Section A: Pure Mathematics}\par
\[1 Write down the general term in the expansion in powers of x of (1 x^6) -2 .\]
-\par
\textbf{(i)} Find the coefficient of x24 in the expansion in powers of x of\par
6 -2 3 -1\par
(1 x ) (1 x ) .\par
- -\par
Obtain also, and simplify, formulae for the coefficient of xn in the different cases that\par
arise.\par
\textbf{(ii)} Show that the coefficient of x24 in the expansion in powers of x of\par
6 -2 3 -1 -1\par
(1 x ) (1 x ) (1 x)\par
- - -\par
\[is 55, and find the coefficients of x^25 and x^66.\]
2 If p(x) and q(x) are polynomials of degree m and n, respectively, what is the degree of p(q(x))?\par
\textbf{(i)} The polynomial p(x) satisfies\par
\[p(p(p(x))) 3p(x) = 2x\]
- -\par
for all x. Explain carefully why p(x) must be of degree 1, and find all polynomials that\par
satisfy this equation.\par
\textbf{(ii)} Find all polynomials that satisfy\par
2 4\par
\[2p(p(x)) + 3[p(x)] 4p(x) = x\]
-\par
for all x.\par
\end{flushleft}
\subsection*{Page 2}
\begin{flushleft}
\[\pi\]
\bigskip\textbf{Question 3.} Show that, for any function f (for which the integrals exist),\par
\infty \infty\par
1 1\par
\[f x + 1 + x^2 dx = 1 + f(t) dt .\]
\[\leq0\]
\[\geq ( )\]
\[2 \leq1 \int t^2 \sum\]
Hence evaluate\par
\infty\par
dx ,\par
\[\leq0 2x2 + 1 + 2x\sqrt\{x^2\} + 1\]
\[and, using the substitution x = tan θ,\]
\[2 \pi 1\]
dθ .\par
\[\leq0 (1 + sin θ)3\]
\bigskip\textbf{Question 4.} In this question, you may assume that the infinite series\par
\[x^2 x^3 x^4 xn\]
\[ln(1 + x) = x + + + (\]
1)n+1\par
+\par
- 2 3 - 4 \cdots - n \cdots\par
\[is valid for x < 1.\]
| |\par
\textbf{(i)} Let n be an integer greater than 1. Show that, for any positive integer k,\par
1 1\par
\[< .\]
(k + 1)nk+1 knk\par
1 1\par
\[Hence show that ln 1 + < . Deduce that\]
\[\int n \sum n\]
n\par
\[1 + < e .\]
\[\int n \sum\]
1 2y + 1 1\par
\textbf{(ii)} Show, using an expansion in powers of , that ln > for y > .\par
\[y \int 2y 1 \sum y 2\]
-\par
Deduce that, for any positive integer n,\par
n+ 1\par
1 2\par
\[e < 1 + .\]
\[\int n \sum\]
\textbf{(iii)} Use parts (i) and (ii) to show that as n\par
\to \infty\par
n\par
1 + e .\par
\[\int n \sum \to\]
\end{flushleft}
\subsection*{Page 3}
\begin{flushleft}
\sigma\par
\bigskip\textbf{Question 5.} (i) Sketch the curve y = f(x), where\par
\[f(x) = (x = a 1),\]
(x a)2 1 ) ±\par
- -\par
and a is a constant.\par
\textbf{(ii)} The function g(x) is defined by\par
\[g(x) = (x = a 1, x = b 1),\]
(x a)2 1 (x b)2 1 ) ± ) ±\par
- - - -\par
\[\geq )\geq )\]
\[where a and b are constants, and b > a. Sketch the curves y = g(x) in the two cases\]
\[b > a + 2 and b = a + 2, finding the values of x at the stationary points.\]
\bigskip\textbf{Question 6.} A cyclic quadrilateral ABCD has sides AB, BC, CD and DA of lengths a, b, c and d,\par
respectively. The area of the quadrilateral is Q, and angle DAB is θ.\par
Find an expression for cos θ in terms of a, b, c and d, and an expression for sin θ in terms of\par
a, b, c, d and Q. Hence show that\par
2 2 2 2 2 2 2\par
\[16Q = 4(ad + bc) (a + d b c ) ,\]
- - -\par
and deduce that\par
\[Q = (s a)(s b)(s c)(s d) ,\]
- - - -\par
\[where s = (a + b + c + d).\]
Deduce a formula for the area of a triangle with sides of length a, b and c.\par
\bigskip\textbf{Question 7.} Three distinct points, X1, X2 and X3, with position vectors x1, x2 and x3 respectively, lie\par
on a circle of radius 1 with its centre at the origin O. The point G has position vector\par
\[(x^1 + x^2 + x^3). The line through X^1 and G meets the circle again at the point Y^1 and the\]
\[points Y^2 and Y^3 are defined correspondingly.\]
\[Given that -G-Y\to1 = λ1-G-X\to1, where λ1 is a positive scalar, show that\]
-\par
\[-O-Y\to1 =\]
(1\par
-\par
\[2λ1)x^1 + (1 + λ1)(x^2 + x^3)\]
\[\geq )\]
and hence that\par
3 α β γ\par
\[λ1 = - - - ,\]
3 + α 2β 2γ\par
- -\par
\[where α = x^2 . x^3, β = x^3 . x^1 and γ = x^1 . x^2.\]
\[GX1 GX2 GX3\]
\[Deduce that + + = 3 .\]
\[GY1 GY2 GY3\]
\end{flushleft}
\subsection*{Page 4}
\begin{flushleft}
\omega\par
\bigskip\textbf{Question 8.} The positive numbers α, β and q satisfy β α > q. Show that\par
-\par
\[α2 + β2 q^2\]
\[- 2 > 0 .\]
αβ -\par
\[The sequence u^0, u^1, . . . is defined by u^0 = α, u^1 = β and\]
\[u^2 q^2\]
\[u n+1 = n - (n \leq 1),\]
u n-1\par
where α, β and q are given positive numbers (and α and β are such that no term in the\par
\[sequence is zero). Prove that u n (u n + u n+2) = u n+1(u n-1 + u n+1) . Prove also that\]
\[u n+1 pu n + u n-1 = 0\]
-\par
for some number p which you should express in terms of α, β and q.\par
\[Hence, or otherwise, show that if β > α + q, the sequence is strictly increasing (that is,\]
\[u n+1 u n > 0 for all n). Comment on the case β = α + q.\]
-\par
\end{flushleft}
\subsection*{Page 5}
\begin{flushleft}
)\par
\bigskip\textbf{Section B: Mechanics}\par
\bigskip\textbf{Question 9.} A tennis ball is projected from a height of 2h above horizontal ground with speed u and at an\par
angle of α below the horizontal. It travels in a plane perpendicular to a vertical net of height\par
h which is a horizontal distance of a from the point of projection. Given that the ball passes\par
over the net, show that\par
1 2(h a tan α)\par
\[< - .\]
\[u^2 ga2 sec2 α\]
The ball lands before it has travelled a horizontal distance of b from the point of projection.\par
Show that\par
bg\par
\[u^2 sin 2 α + 4gh < + u sin α .\]
( u cos α\par
Hence show that\par
\[h(b^2 2a2)\]
\[tan α < - .\]
ab(b a)\par
-\par
\bigskip\textbf{Question 10.} A hollow circular cylinder of internal radius r is held fixed with its axis horizontal. A uniform\par
rod of length 2a (where a < r) rests in equilibrium inside the cylinder inclined at an angle of\par
θ to the horizontal, where θ = 0. The vertical plane containing the rod is perpendicular to\par
)\par
the axis of the cylinder. The coefficient of friction between the cylinder and each end of the\par
\[rod is μ, where μ > 0.\]
Show that, if the rod is on the point of slipping, then the normal reactions R1 and R2 of the\par
lower and higher ends of the rod, respectively, on the cylinder are related by\par
\[μ(R^1 + R^2) = (R^1 R^2) tan φ\]
-\par
where φ is the angle between the rod and the radius to an end of the rod.\par
Show further that\par
\[μr2\]
\[tan θ = .\]
\[r^2 a^2(1 + μ2)\]
-\par
\[Deduce that λ < φ, where tan λ = μ.\]
\end{flushleft}
\subsection*{Page 6}
\begin{flushleft}
(\par
\bigskip\textbf{Question 11.} A small block of mass km is initially at rest on a smooth horizontal surface. Particles P1, P2,\par
P3, . . . are fired, in order, along the surface from a fixed point towards the block. The mass\par
of the ith particle is im (i = 1, 2, . . .) and the speed at which it is fired is u/i . Each particle\par
that collides with the block is embedded in it. Show that, if the nth particle collides with the\par
block, the speed of the block after the collision is\par
2nu\par
.\par
2k + n(n + 1)\par
In the case 2k = N (N + 1), where N is a positive integer, determine the number of collisions\par
that occur. Show that the total kinetic energy lost in all the collisions is\par
N+1\par
1 2\par
mu .\par
\[2 \int n \sum\]
\[n)=2\]
\end{flushleft}
\subsection*{Page 7}
\begin{flushleft}
)\par
\bigskip\textbf{Section C: Probability and Statistics}\par
\bigskip\textbf{Question 12.} A modern villa has complicated lighting controls. In order for the light in the swimming pool\par
to be on, a particular switch in the hallway must be on and a particular switch in the kitchen\par
must be on. There are four identical switches in the hallway and four identical switches in the\par
kitchen. Guests cannot tell whether the switches are on or off, or what they control.\par
Each Monday morning a guest arrives, and the switches in the hallway are either all on or all\par
off. The probability that they are all on is p and the probability that they are all off is 1 p.\par
1 -\par
The switches in the kitchen are each on or off, independently, with probability .\par
\textbf{(i)} On the first Monday, a guest presses one switch in the hallway at random and one switch\par
in the kitchen at random. Find the probability that the swimming pool light is on at the\par
end of this process. Show that the probability that the guest has pressed the swimming\par
pool light switch in the hallway, given that the light is on at the end of the process,\par
1 p\par
is - .\par
1 + 2p\par
\textbf{(ii)} On each of seven Mondays, guests go through the above process independently of each\par
other, and each time the swimming pool light is found to be on at the end of the process.\par
Given that the most likely number of days on which the swimming pool light switch in\par
1 5\par
\[the hallway was pressed is 3, show that < p < .\]
4 14\par
\infty\par
\bigskip\textbf{Question 13.} In this question, you may assume that e\par
\[-x^2/2\]
\[dx =\]
\[\pi .\]
\[\leq0 (\]
The number of supermarkets situated in any given region can be modelled by a Poisson random\par
variable, where the mean is k times the area of the given region. Find the probability that\par
there are no supermarkets within a circle of radius y.\par
The random variable Y denotes the distance between a randomly chosen point in the region\par
\[and the nearest supermarket. Write down P(Y < y) and hence show that the probability\]
\[density function of Y is 2\piyke\]
\[-\piky2\]
\[for y \leq 0.\]
\[4 \pi\]
\[Find E(Y ) and show that Var (Y ) = - .\]
\[4\pik\]
\end{flushleft}
\end{document}
